\chapter[Lezione 10 --- 20 ottobre 2026]{Lezione 10: applicazioni inverse e monotonia}
\label{chap:lezione-10}

\section{Dominio, immagine e grafico dell'applicazione inversa}
\label{sec:lez10-proprieta-inversa}

La \cref{def:lez09-inversa} permette di precisare il rapporto
tra dominio e immagine di un'applicazione iniettiva e della sua inversa.

\begin{proposition}[Dominio e immagine dell'applicazione inversa]
\label{prop:lez10-dominio-immagine-inversa}
Siano \(A,B\) insiemi non vuoti e
\(\varphi\colon A\longrightarrow B\) un'applicazione iniettiva.
Allora
\[
  D(\varphi^{-1})=\operatorname{Im}(\varphi),
  \qquad
  \operatorname{Im}(\varphi^{-1})=D(\varphi).
\]
\end{proposition}

\begin{proof}
\emph{Prima uguaglianza.}
Segue direttamente dalla \cref{def:lez09-inversa}.

\emph{Seconda uguaglianza.}
Poiché \(C(\varphi^{-1})=A=D(\varphi)\),
per la \cref{def:lez06-immagine} si ha
\[
  \operatorname{Im}(\varphi^{-1})\subseteq D(\varphi).
\]
Viceversa, sia \(a_0\in D(\varphi)\).
Ponendo \(b_0:=\varphi(a_0)\), si ottiene
\(b_0\in\operatorname{Im}(\varphi)=D(\varphi^{-1})\).
Dalla \cref{def:lez09-inversa} segue
\[
  \varphi^{-1}(b_0)=a_0,
\]
quindi \(a_0\in\operatorname{Im}(\varphi^{-1})\).
Pertanto \(D(\varphi)\subseteq\operatorname{Im}(\varphi^{-1})\),
da cui la seconda uguaglianza.
\end{proof}

\begin{example}[Esponenziale e logaritmo]
\label{ex:lez10-dominio-immagine-logaritmo}
L'applicazione
\[
  \varphi\colon\R\longrightarrow(0,+\infty),
  \qquad x\longmapsto e^x,
\]
è invertibile e la sua inversa è
\[
  \varphi^{-1}\colon(0,+\infty)\longrightarrow\R,
  \qquad y\longmapsto\log y.
\]
In accordo con la \cref{prop:lez10-dominio-immagine-inversa},
\[
  D(\varphi^{-1})=(0,+\infty),
  \qquad \operatorname{Im}(\varphi^{-1})=\R.
\]
\end{example}

Se \(A,B\subseteq\R\), il grafico dell'inversa si ottiene
scambiando le coordinate dei punti del grafico originario.
Infatti, per la \cref{def:lez07-grafico},
\[
\begin{aligned}
  G(\varphi)
    &=\{(a,\varphi(a)):a\in A\},\\
  G(\varphi^{-1})
    &=\{(b,\varphi^{-1}(b)):b\in\operatorname{Im}(\varphi)\}\\
    &=\{(\varphi(a),a):a\in A\}.
\end{aligned}
\]

\clearpage

I due grafici sono quindi simmetrici rispetto alla retta
\(y=x\), rappresentata nella \cref{fig:lez10-grafico-inversa}.

\begin{figure}[htbp]
  \centering
  \begin{tikzpicture}[unipd/diagram,x=1.12cm,y=1.12cm]
  \clip (-2.35,-2.35) rectangle (4.35,4.25);
  \draw[->,unipd-muted] (-2.25,0) -- (4.23,0)
    node[below right] {\(x\)};
  \draw[->,unipd-muted] (0,-2.25) -- (0,4.12)
    node[above left] {\(y\)};
  \draw[unipd-muted,densely dashed,thick] (-2.15,-2.15) -- (4.05,4.05)
    node[above left] {\(y=x\)};
  \draw[very thick,unipd-primary,domain=-2.2:1.39,samples=110]
    plot (\x,{exp(\x)});
  \draw[very thick,unipd-accent,domain=.11:4.13,samples=110]
    plot (\x,{ln(\x)});
  \fill[unipd-primary] (0,1) circle (1.9pt);
  \fill[unipd-accent] (1,0) circle (1.9pt);
  \node[left=4pt] at (0,1) {\((0,1)\)};
  \node[below=5pt] at (1,0) {\((1,0)\)};
  \node[text=unipd-primary,above left] at (1.32,3.74) {\(y=e^x\)};
  \node[text=unipd-accent,above left] at (3.65,1.28) {\(y=\log x\)};
\end{tikzpicture}

  \caption{I grafici di \(e^x\) e \(\log x\) sono simmetrici
  rispetto alla retta \(y=x\).}
  \label{fig:lez10-grafico-inversa}
\end{figure}

\section{Monotonia delle funzioni}
\label{sec:lez10-monotonia}

Siano \(A\subseteq\R\) non vuoto,
\(\varphi\colon A\longrightarrow\R\) e
\(S\subseteq A\) non vuoto. Le definizioni seguenti
si applicano, in particolare, quando \(S=I\) è un intervallo,
nel senso della \cref{def:intervallo-retta-reale}.

\begin{definition}[Crescenza, decrescenza e monotonia]
\label{def:lez10-monotonia}
Per ogni \(x_1,x_2\in S\) con \(x_1<x_2\),
si dice che \(\varphi\) è su \(S\):
\begin{enumerate}
  \item \emph{strettamente crescente} se
  \(\varphi(x_1)<\varphi(x_2)\);
  \item \emph{debolmente crescente} se
  \(\varphi(x_1)\le\varphi(x_2)\);
  \item \emph{strettamente decrescente} se
  \(\varphi(x_1)>\varphi(x_2)\);
  \item \emph{debolmente decrescente} se
  \(\varphi(x_1)\ge\varphi(x_2)\).
\end{enumerate}
Si dice che \(\varphi\) è \emph{strettamente monotona} su \(S\)
se è strettamente crescente oppure strettamente decrescente,
e \emph{debolmente monotona} su \(S\) se è debolmente crescente
oppure debolmente decrescente.
\end{definition}

\begin{example}[Un'unione che non è un intervallo]
\label{ex:lez10-unione-non-intervallo}
L'insieme
\[
  J=(1,2)\cup(2,4)
\]
non è un intervallo, benché sia unione di due intervalli.
Infatti, scegliendo
\[
  x_1=\frac32\in J,
  \qquad x_2=3\in J,
\]
si ha \(x_1<2<x_2\) ma \(2\notin J\),
in contrasto con la \cref{def:intervallo-retta-reale}.
\end{example}

\clearpage

\begin{proposition}[Monotonia della funzione opposta]
\label{prop:lez10-monotonia-opposta}
Siano \(I\subseteq\R\) un intervallo non vuoto e
\(\varphi\colon I\longrightarrow\R\).

Se \(\varphi\) è debolmente crescente su \(I\),
allora \(-\varphi\) è debolmente decrescente su \(I\).
Se \(\varphi\) è debolmente decrescente su \(I\),
allora \(-\varphi\) è debolmente crescente su \(I\).
Le stesse proprietà valgono per la monotonia stretta.
\end{proposition}

\begin{proof}
Siano \(x_1,x_2\in I\) con \(x_1<x_2\).

\emph{Caso crescente.}
Se \(\varphi(x_1)\le\varphi(x_2)\), cambiando
segno a entrambi i membri si inverte il verso
della disuguaglianza e si ottiene
\[
  (-\varphi)(x_1)=-\varphi(x_1)
  \ge-\varphi(x_2)=(-\varphi)(x_2).
\]
Pertanto \(-\varphi\) è debolmente decrescente.

\emph{Caso decrescente.}
Se \(\varphi(x_1)\ge\varphi(x_2)\),
con lo stesso procedimento si ottiene
\[
  (-\varphi)(x_1)=-\varphi(x_1)
  \le-\varphi(x_2)=(-\varphi)(x_2).
\]
Pertanto \(-\varphi\) è debolmente crescente.
Se le disuguaglianze iniziali sono strette,
restano strette anche quelle ottenute moltiplicando per \(-1\).
\end{proof}

\begin{proposition}[Monotonia del prodotto di funzioni non negative]
\label{prop:lez10-monotonia-prodotto}
Siano \(I\subseteq\R\) un intervallo non vuoto e
\(\varphi,\psi\colon I\longrightarrow\R\), con
\[
  \operatorname{Im}(\varphi)\subseteq[0,+\infty),
  \qquad
  \operatorname{Im}(\psi)\subseteq[0,+\infty).
\]
Se \(\varphi\) e \(\psi\) sono entrambe debolmente crescenti
su \(I\), allora il prodotto \(\varphi\psi\) è debolmente
crescente su \(I\).

Se \(\varphi\) e \(\psi\) sono entrambe debolmente decrescenti
su \(I\), allora il prodotto \(\varphi\psi\) è debolmente
decrescente su \(I\).
\end{proposition}

\begin{proof}
Siano \(x_1,x_2\in I\) con \(x_1<x_2\).

\emph{Caso crescente.}
Per ipotesi,
\[
  0\le\varphi(x_1)\le\varphi(x_2),
  \qquad 0\le\psi(x_1)\le\psi(x_2).
\]
Ne segue che
\[
  \varphi(x_1)\psi(x_1)
  \le\varphi(x_2)\psi(x_1)
  \le\varphi(x_2)\psi(x_2),
\]
quindi \(\varphi\psi\) è debolmente crescente.

\emph{Caso decrescente.}
Per ipotesi,
\[
  \varphi(x_1)\ge\varphi(x_2)\ge0,
  \qquad \psi(x_1)\ge\psi(x_2)\ge0.
\]
Ne segue che
\[
  \varphi(x_1)\psi(x_1)
  \ge\varphi(x_2)\psi(x_1)
  \ge\varphi(x_2)\psi(x_2),
\]
quindi \(\varphi\psi\) è debolmente decrescente.
\end{proof}

Quando una funzione è crescente e l'altra decrescente,
non si può dedurre, in generale, la monotonia del prodotto.
Ad esempio, su \(I=[0,1]\) le funzioni non negative
\(\varphi(x)=x\) e \(\psi(x)=1-x\) sono rispettivamente
crescenti e decrescenti, ma
\[
  (\varphi\psi)(0)=(\varphi\psi)(1)=0,
  \qquad (\varphi\psi)\!\left(\frac12\right)=\frac14,
\]
per cui il prodotto non è monotono su \(I\).

\begin{proposition}[Monotonia stretta e applicazione inversa]
\label{prop:lez10-monotonia-inversa}
Siano \(I\subseteq\R\) un intervallo non vuoto e
\(\varphi\colon I\longrightarrow\R\) strettamente monotona su \(I\).
Allora:
\begin{enumerate}
  \item \(\varphi\) è iniettiva;
  \item l'applicazione inversa
  \(\varphi^{-1}\colon\operatorname{Im}(\varphi)\longrightarrow I\)
  è strettamente monotona su \(\operatorname{Im}(\varphi)\).
  In particolare, se \(\varphi\) è strettamente crescente,
  anche \(\varphi^{-1}\) è strettamente crescente;
  se \(\varphi\) è strettamente decrescente,
  anche \(\varphi^{-1}\) è strettamente decrescente.
\end{enumerate}
\end{proposition}

\begin{proof}[Dimostrazione del primo punto]
Siano \(x_1,x_2\in I\) distinti.
Se \(x_1<x_2\), la stretta monotonia implica
\(\varphi(x_1)<\varphi(x_2)\), nel caso crescente, e
\(\varphi(x_1)>\varphi(x_2)\), nel caso decrescente.
Se \(x_2<x_1\) si procede analogamente.
In entrambi i casi \(\varphi(x_1)\ne\varphi(x_2)\),
dunque \(\varphi\) è iniettiva per la
\cref{def:lez06-iniettiva}.
\end{proof}

La seconda affermazione della
\cref{prop:lez10-monotonia-inversa} è enunciata senza dimostrazione.

\begin{example}[Monotonia del logaritmo]
\label{ex:lez10-logaritmo-monotono}
Per ogni \(a>1\), la funzione
\[
  \varphi\colon\R\longrightarrow(0,+\infty),
  \qquad x\longmapsto a^x,
\]
è strettamente crescente su \(\R\).
La sua inversa
\[
  \varphi^{-1}\colon(0,+\infty)\longrightarrow\R,
  \qquad x\longmapsto\log_a x,
\]
è quindi strettamente crescente su \((0,+\infty)\)
per la \cref{prop:lez10-monotonia-inversa}.
\end{example}
